\thesisTitle{EMPIRICAL IMPACT OF GRAPH OPTIMIZATIONS (OPERATOR FUSION) ON NPU INFERENCE}
\thesisTitleTR{NPU ÜZERİNDE ÇIKARIMDA GRAFİK OPTİMİZASYONLARININ (OPERATÖR FÜZYONU) AMPİRİK ETKİSİ}
%\thesisYear{2026}
%\thesisMonth{April}\thesisMonthTR{Nisan}
\thesisDefenseDay{12}  % day of defense, 1, 2, ..., 31
\thesisDegree{MSc}  %% PhD || MSc
\thesisDegreeTR{Yüksek Lisans} %% Doktora || Yüksek Lisans
\thesisProgram{Software Engineering}
\thesisProgramTR{Yazılım Mühendisliği}
\thesisAuthor{Yusuf Talha ARABACI}
\thesisAdvisorTitle{Dr. Öğr. Üyesi}
\thesisAdvisorTitleTR{Dr. Öğr. Üyesi}
\thesisAdvisor{Emrullah DEMİRAL}
\thesisAdvisorDepartment{Department of Software Engineering}

% Template customization for course project (overrides defaults in the class)
\thesisDocTypeEN{COURSE PROJECT REPORT}
\thesisDocTypeTR{DERS PROJESİ RAPORU}
\thesisAdvisorLabelEN{Project Advisor}
\thesisAdvisorLabelTR{Proje Danışmanı}
\thesisPreparedAsEN{Course Project Report}
\thesisPreparedAsTR{Ders Projesi Raporu}
\thesisCourseNameEN{Advanced Algorithm Analysis -- Course Project}
\thesisCourseNameTR{İleri Algoritma Analizi -- Ders Projesi}

\thesisChairman{(Jüri Başkanı)}
\thesisMemberTWO{(Üye 1)}
\thesisMemberTHREE{(Üye 2)}
\thesisMemberFOUR{}
\thesisMemberFIVE{}

\thesisPages{29}

\thesisAbstract{ 
This report presents an empirical study prepared as a course project for \textbf{Advanced Algorithm Analysis}. It investigates how graph-level optimizations---with a focus on operator fusion---affect end-to-end inference latency on NPU-class accelerators.
The experimental pipeline is built on ONNX Runtime and compares baseline execution (optimizations disabled) against an optimized configuration (extended graph optimizations enabled) \cite{onnxruntime}.
We run three representative ONNX models (MobileNetV2, ResNet50, and BERT-tiny) under a provider fallback chain that prioritizes the OpenVINO Execution Provider \cite{openvino_ep}.
We collect latency statistics across repeated runs and export ONNX Runtime profiling traces to quantify operator-level time breakdowns and graph node reductions.
Results show that node count can be reduced substantially (\(27\%\) to \(68\%\)), yet the observed latency impact is model-dependent: some models regress while others improve.
We discuss these outcomes in the context of the memory wall \cite{memorywall} and accelerator runtime overheads, emphasizing that fusion and graph rewrites do not guarantee speedups without considering EP-specific compilation and kernel selection.
}

\thesisKeywords{ONNX Runtime, graph optimization, operator fusion, NPU, OpenVINO, profiling, performance benchmarking.}

\thesisOzet{ 
Bu rapor, \textbf{İleri Algoritma Analizi} kapsamında hazırlanmış bir ders projesi olarak, operatör füzyonuna odaklanarak grafik-seviyesi optimizasyonların NPU sınıfı hızlandırıcılarda uçtan uca çıkarım gecikmesine etkisini ampirik olarak inceler.
Deneysel altyapı ONNX Runtime üzerinde kurulmuş olup optimizasyonların kapalı olduğu bir temel durum ile genişletilmiş optimizasyonların açık olduğu bir optimize durum karşılaştırılmıştır \cite{onnxruntime}.
MobileNetV2, ResNet50 ve BERT-tiny olmak üzere üç ONNX modelinde, öncelikli olarak OpenVINO Execution Provider kullanan ve gerektiğinde CPU'ya düşen bir sağlayıcı zinciriyle çalıştırma yapılmıştır \cite{openvino_ep}.
Tekrarlı ölçümlerle gecikme istatistikleri çıkarılmış, ONNX Runtime profiling izleri dışa aktarılmış ve operatör-seviyesi süre kırılımları ile grafik düğüm azalımı nicelleştirilmiştir.
Sonuçlar düğüm sayısında belirgin azalmalar (\(\%27\) -- \(\%68\)) elde edilebildiğini, ancak gecikme etkisinin modele bağlı olduğunu göstermektedir.
Elde edilen bulgular, bellek duvarı \cite{memorywall} ve sağlayıcıya özgü derleme/çalıştırma ek yükleri çerçevesinde tartışılmış; füzyonun her koşulda hızlanma garanti etmediği vurgulanmıştır.
}

\thesisKeywordsTR{ONNX Runtime, grafik optimizasyonu, operatör füzyonu, NPU, OpenVINO, profiling, performans kıyaslama.}

\thesisScienceCode{0}

\thesisAcknowledgment{Bu çalışmanın yürütülmesinde desteği olan ders yürütücüsüne ve katkı sağlayan herkese teşekkür ederim.}

\thesisResume{Yusuf Talha ARABACI (yusuftalhaarabaci@hotmail.com) -- Yazılım Mühendisliği Yüksek Lisans öğrencisi. İlgi alanları: derleyici optimizasyonları, grafik optimizasyonları, hızlandırıcılar ve performans analizi.}

\thesisSymbolList{
\thesisSymbol{$S$}{Speedup (baseline/optimized)}
\thesisSymbol{$t$}{Inference latency (ms)}
\thesisSymbol{$N$}{Graph node count}
\thesisSymbol{$AI$}{Arithmetic intensity (FLOP/byte)}
}

\thesisAbbrList{
\thesisAbbr{ORT}{ONNX Runtime}
\thesisAbbr{EP}{Execution Provider}
\thesisAbbr{NPU}{Neural Processing Unit}
\thesisAbbr{CPU}{Central Processing Unit}
\thesisAbbr{OV}{OpenVINO}
\thesisAbbr{FLOP}{Floating-point operation}
}

\makeThesisPreamble
\pagenumbering{arabic}